% ============================================================
%  CHAPITRE : LIMITES ET CONTINUITÉ
%  Mazava Maths — Terminale S (Bacc Série S, Madagascar)
%  Harmonisation : fusion de « Limites de fonctions » et de la
%  partie « Continuité » de l'ancien chapitre Continuité/Dérivabilité.
%  Parti pris : tous les exemples et exercices utilisent les
%  fonctions ln et exp (les puissances n'apparaissent que comme
%  terme de comparaison dans les croissances comparées).
%  Environnements : definition, propriete, theoreme, methode,
%                   attention, astucebac, exemple, exercice,
%                   solution, remarque
% ============================================================

\chapter{Limites et continuité}

\begin{citationchapitre}
	« Approchez, et la foi vous viendra. »\\[0.4em]
	— d'après Jean le Rond d'Alembert
\end{citationchapitre}

\begin{objectifs}
	\begin{itemize}[label=\textbullet]
		\item Déterminer la limite d'une fonction bâtie sur $\ln$ et $\exp$
		en une borne de son ensemble de définition.
		\item Reconnaître une forme indéterminée et la lever à l'aide des
		\textbf{limites usuelles} et des \textbf{croissances comparées}.
		\item Étudier la limite d'une \textbf{fonction composée} ; appliquer les
		théorèmes de \textbf{comparaison} et des \textbf{gendarmes}.
		\item Étudier la \textbf{continuité} d'une fonction en un point et sur un
		intervalle ; prolonger une fonction par continuité.
		\item Utiliser le \textbf{théorème des valeurs intermédiaires} (TVI) et son
		corollaire pour une fonction strictement monotone.
		\item Interpréter limites et continuité : \textbf{asymptotes}, \textbf{image
		d'un intervalle}.
	\end{itemize}
\end{objectifs}


% ================================================================
\section{La notion de limite (rappels)}
% ================================================================

\begin{definition}
	Soit $f$ une fonction et $\ell \in \R$.
	\begin{itemize}[label=\textbullet]
		\item $\displaystyle\lim_{x\to+\infty}f(x)=\ell$ : $f(x)$ est aussi proche
		de $\ell$ que l'on veut dès que $x$ est assez grand.
		\item $\displaystyle\lim_{x\to+\infty}f(x)=+\infty$ : $f(x)$ dépasse tout
		réel dès que $x$ est assez grand.
		\item $\displaystyle\lim_{x\to a}f(x)=\ell$ : $f(x)$ se rapproche de $\ell$
		quand $x$ se rapproche de $a$ (sans forcément l'atteindre).
		\item $\displaystyle\lim_{x\to a}f(x)=\pm\infty$ : $f(x)$ devient infini
		quand $x\to a$. On distingue si besoin $x\to a^-$ et $x\to a^+$.
	\end{itemize}
\end{definition}

\begin{remarque}
	Si $\displaystyle\lim_{x\to a^-}f(x)\neq\lim_{x\to a^+}f(x)$, alors $f$
	\textbf{n'a pas de limite} en $a$.
\end{remarque}

\begin{propriete}{Limites des puissances (pour comparer)}{prop_lim_puissances}
	\begin{center}
	\renewcommand{\arraystretch}{1.5}
	\begin{tabular}{>{\centering\arraybackslash}p{2.4cm}
	                >{\centering\arraybackslash}p{2.2cm}
	                >{\centering\arraybackslash}p{2.2cm}}
	\hline
	\textbf{Fonction} & \textbf{Variable} & \textbf{Limite} \\
	\hline
	$x^n$ ($n\geq 1$)      & $x\to+\infty$   & $+\infty$ \\
	$\dfrac{1}{x^n}$       & $x\to\pm\infty$ & $0$ \\
	$\dfrac{1}{x}$         & $x\to 0^+$      & $+\infty$ \\
	$\dfrac{1}{x}$         & $x\to 0^-$      & $-\infty$ \\
	\hline
	\end{tabular}
	\end{center}
\end{propriete}


% ================================================================
\section{Limites de référence : \texorpdfstring{$\ln$}{ln} et \texorpdfstring{$\exp$}{exp}}
% ================================================================

\begin{propriete}{Limites aux bornes}{prop_lim_bornes}
	\begin{center}
	\renewcommand{\arraystretch}{1.7}
	\begin{tabular}{>{\centering\arraybackslash}p{4.6cm}
	                >{\centering\arraybackslash}p{2.2cm}}
	\hline
	\textbf{Limite} & \textbf{Résultat} \\
	\hline
	$\displaystyle\lim_{x\to 0^+}\ln x$     & $-\infty$ \\
	$\displaystyle\lim_{x\to +\infty}\ln x$ & $+\infty$ \\
	\hline
	$\displaystyle\lim_{x\to -\infty}\e^x$  & $0^+$ \\
	$\displaystyle\lim_{x\to +\infty}\e^x$  & $+\infty$ \\
	\hline
	\end{tabular}
	\end{center}
\end{propriete}

\begin{center}
\begin{tikzpicture}[>=Stealth, scale=0.9]
	\draw[->] (-3.3,0) -- (4.4,0) node[right,font=\small] {$x$};
	\draw[->] (0,-3.3) -- (0,3.7) node[above,font=\small] {$y$};
	\draw[gray!45,dashed] (-3.1,-3.1) -- (3.4,3.4);
	\node[font=\scriptsize,gray,anchor=west] at (3.15,3.35) {$y=x$};
	\draw[thick,domain=-3.2:1.1,samples=100] plot (\x,{exp(\x)});
	\draw[thick,domain=0.05:4.2,samples=140] plot (\x,{ln(\x)});
	\draw (1,-0.09) -- (1,0.09) node[below,font=\scriptsize] {$1$};
	\draw (-0.09,1) -- (0.09,1) node[left,font=\scriptsize] {$1$};
	\node[font=\small,anchor=west] at (1.35,2.55) {$y=\e^x$};
	\node[font=\small,anchor=west] at (2.9,1.55) {$y=\ln x$};
\end{tikzpicture}
\end{center}

\begin{propriete}{Croissances comparées}{prop_croissances_comparees}
	Pour tout entier $n\geq 1$ :
	\begin{center}
	\renewcommand{\arraystretch}{2.0}
	\begin{tabular}{>{\raggedright\arraybackslash}p{6.0cm}
	                >{\centering\arraybackslash}p{2.0cm}}
	\hline
	$\displaystyle\lim_{x\to+\infty}\frac{\ln x}{x^n}$              & $0$ \\
	$\displaystyle\lim_{x\to 0^+}x^n\ln x$                          & $0$ \\
	$\displaystyle\lim_{x\to+\infty}\frac{\e^x}{x^n}$               & $+\infty$ \\
	$\displaystyle\lim_{x\to-\infty}x^n\,\e^x$                      & $0$ \\
	\hline
	\end{tabular}
	\end{center}
\end{propriete}

\begin{propriete}{Limites usuelles (taux de variation en $0$ et en $1$)}{prop_lim_usuelles}
	\begin{center}
	\renewcommand{\arraystretch}{2.0}
	\setlength{\tabcolsep}{4pt}
	\begin{tabular}{>{\raggedright\arraybackslash}p{4.3cm}
	                >{\raggedleft\arraybackslash}p{4.3cm}}
	\hline
	$\displaystyle\lim_{x\to 0}\frac{\e^x-1}{x}=1$ &
	$\bigl(\e^x\bigr)'_{x=0}=1$ \\
	$\displaystyle\lim_{x\to 0}\frac{\ln(1+x)}{x}=1$ &
	$\bigl(\ln(1+x)\bigr)'_{x=0}=1$ \\
	$\displaystyle\lim_{x\to 1}\frac{\ln x}{x-1}=1$ &
	$\bigl(\ln x\bigr)'_{x=1}=1$ \\
	\hline
	\end{tabular}
	\end{center}
\end{propriete}

\begin{astucebac}
	\textbf{Hiérarchie à l'infini :}
	\[
	\e^x \;\gg\; x^n \;\gg\; \ln x \;\gg\; 1 .
	\]
	« L'exponentielle écrase toute puissance ; le logarithme se fait écraser. »
	Dans un quotient, on met en facteur \textbf{le terme le plus fort}
	(le plus à gauche) au numérateur \emph{et} au dénominateur, puis on simplifie.
\end{astucebac}


% ================================================================
\section{Opérations et formes indéterminées}
% ================================================================

\begin{propriete}{Règles opératoires}{prop_operations}
	Si $\displaystyle\lim_{x\to a}f=\ell$ et $\displaystyle\lim_{x\to a}g=m$
	($\ell,m\in\R$ ou $\pm\infty$) :
	\begin{center}
	\renewcommand{\arraystretch}{1.7}
	\begin{tabular}{>{\centering\arraybackslash}p{1.6cm}
	                >{\centering\arraybackslash}p{3.0cm}
	                >{\centering\arraybackslash}p{3.4cm}}
	\hline
	\textbf{Opé.} & \textbf{Résultat} & \textbf{Forme indéterminée} \\
	\hline
	$f+g$ & $\ell+m$ & $(+\infty)+(-\infty)$ \\
	$f\times g$ & $\ell\times m$ & $0\times\infty$ \\
	$\dfrac{f}{g}$ & $\dfrac{\ell}{m}$ ($m\neq 0$) &
	$\dfrac{0}{0}$ \quad ; \quad $\dfrac{\infty}{\infty}$ \\
	\hline
	\end{tabular}
	\end{center}
\end{propriete}

\begin{attention}
	Les quatre \textbf{formes indéterminées} (FI) sont
	\[
	\frac{0}{0},\qquad \frac{\infty}{\infty},\qquad
	\infty-\infty,\qquad 0\times\infty .
	\]
	Les formes $1^{\infty}$, $\infty^{0}$, $0^{0}$ (rencontrées avec
	$u(x)^{v(x)}=\e^{v(x)\ln u(x)}$) se ramènent à $0\times\infty$ dans
	l'exposant.
\end{attention}


% ================================================================
\section{Lever une indétermination avec les limites usuelles}
% ================================================================

% --------------------------------------------------
\subsection{Mettre en facteur le terme dominant (croissances comparées)}
% --------------------------------------------------

\begin{methode}
	Pour une FI $\infty-\infty$ ou $\dfrac{\infty}{\infty}$ mêlant
	$\e^x$, $x^n$ et $\ln x$ :
	\begin{enumerate}
		\item repérer le terme le plus fort d'après la hiérarchie
		$\e^x\gg x^n\gg\ln x$ ;
		\item le mettre en facteur ;
		\item conclure avec les croissances comparées
		(Propriété~\ref{prop:prop_croissances_comparees}).
	\end{enumerate}
\end{methode}

\begin{exemple}
	$\displaystyle L=\lim_{x\to+\infty}\bigl(\e^x-x^2\bigr)$ (FI $\infty-\infty$).
	\[
	\e^x-x^2=\e^x\!\left(1-\frac{x^2}{\e^x}\right)\!,\qquad
	\frac{x^2}{\e^x}\to 0,\qquad\text{donc } L=+\infty .
	\]
\end{exemple}

\begin{exemple}
	$\displaystyle L=\lim_{x\to+\infty}\frac{2x+1}{\e^x+x}$ (FI $\frac{\infty}{\infty}$).
	On divise haut et bas par $\e^x$ :
	\[
	\frac{2x+1}{\e^x+x}
	=\frac{\dfrac{2x}{\e^x}+\dfrac{1}{\e^x}}{\,1+\dfrac{x}{\e^x}\,}
	\xrightarrow[x\to+\infty]{}\frac{0+0}{1+0}=0 .
	\]
\end{exemple}

\begin{exemple}
	$\displaystyle L=\lim_{x\to+\infty}\bigl(x-\ln x\bigr)$ (FI $\infty-\infty$).
	\[
	x-\ln x=x\!\left(1-\frac{\ln x}{x}\right)\!,\qquad
	\frac{\ln x}{x}\to 0,\qquad\text{donc } L=+\infty .
	\]
\end{exemple}

% --------------------------------------------------
\subsection{Diviser numérateur et dénominateur par le terme dominant}
% --------------------------------------------------

\begin{exemple}
	$\displaystyle L=\lim_{x\to+\infty}\frac{\e^x+2x}{3\e^x-1}$.
	On divise par $\e^x$ :
	\[
	\frac{\e^x+2x}{3\e^x-1}
	=\frac{1+\dfrac{2x}{\e^x}}{\,3-\dfrac{1}{\e^x}\,}
	\xrightarrow[x\to+\infty]{}\frac{1+0}{3-0}=\frac{1}{3}.
	\]
\end{exemple}

\begin{exemple}
	$\displaystyle L=\lim_{x\to-\infty}\frac{\e^x+2}{\e^x-3}$.
	Ici $\e^x\to 0$, il n'y a \textbf{pas} d'indétermination :
	$L=\dfrac{0+2}{0-3}=-\dfrac{2}{3}$.
\end{exemple}

% --------------------------------------------------
\subsection{Reconnaître une limite usuelle}
% --------------------------------------------------

\begin{methode}
	Devant $\dfrac{\e^{\,\square}-1}{\square}$, $\dfrac{\ln(1+\square)}{\square}$
	ou $\dfrac{\ln x}{x-1}$, on fait apparaître exactement la forme usuelle,
	quitte à multiplier/diviser par une constante ou à poser une nouvelle
	variable.
\end{methode}

\begin{exemple}
	$\displaystyle L=\lim_{x\to 0}\frac{\e^{3x}-1}{x}$.
	On écrit
	$\dfrac{\e^{3x}-1}{x}=3\cdot\dfrac{\e^{3x}-1}{3x}$.
	En posant $u=3x\to 0$ : $\dfrac{\e^{u}-1}{u}\to 1$, donc $L=3$.
\end{exemple}

\begin{exemple}
	$\displaystyle L=\lim_{x\to 1}\frac{\ln x}{x^2-1}$ (FI $\frac{0}{0}$).
	\[
	\frac{\ln x}{x^2-1}
	=\frac{\ln x}{x-1}\times\frac{1}{x+1}
	\xrightarrow[x\to 1]{}1\times\frac{1}{2}=\frac{1}{2}.
	\]
\end{exemple}

\begin{exemple}
	$\displaystyle L=\lim_{x\to 0}\frac{\e^x-1}{\ln(1+x)}$ (FI $\frac{0}{0}$).
	\[
	\frac{\e^x-1}{\ln(1+x)}
	=\frac{\dfrac{\e^x-1}{x}}{\dfrac{\ln(1+x)}{x}}
	\xrightarrow[x\to 0]{}\frac{1}{1}=1 .
	\]
\end{exemple}

% --------------------------------------------------
\subsection{Changement de variable}
% --------------------------------------------------

\begin{exemple}
	$\displaystyle L=\lim_{x\to+\infty}x\bigl(\e^{1/x}-1\bigr)$ (FI $\infty\times 0$).
	On pose $X=\dfrac{1}{x}\to 0^+$ :
	\[
	x\bigl(\e^{1/x}-1\bigr)=\frac{\e^{X}-1}{X}\xrightarrow[X\to 0]{}1,
	\qquad\text{donc } L=1 .
	\]
\end{exemple}

\begin{exemple}
	$\displaystyle L=\lim_{x\to+\infty}\frac{\ln x}{\sqrt{x}}$ (FI $\frac{\infty}{\infty}$).
	On pose $X=\sqrt{x}\to+\infty$, donc $\ln x=2\ln X$ :
	\[
	\frac{\ln x}{\sqrt{x}}=\frac{2\ln X}{X}\xrightarrow[X\to+\infty]{}0 .
	\]
\end{exemple}

\begin{exemple}
	$\displaystyle L=\lim_{x\to 0^+}x^{x}$ (FI $0^0$).
	$x^{x}=\e^{x\ln x}$ et $x\ln x\to 0$, donc par composition $L=\e^{0}=1$.
\end{exemple}


% ================================================================
\section{Limite d'une fonction composée}
% ================================================================

\begin{theoreme}[Limite d'une composée]
	Si $\displaystyle\lim_{x\to a}u(x)=b$ et $\displaystyle\lim_{t\to b}g(t)=\ell$
	($a$, $b$, $\ell$ réels ou infinis), alors
	\[
	\lim_{x\to a}g\bigl(u(x)\bigr)=\ell .
	\]
\end{theoreme}

\begin{methode}
	\begin{enumerate}
		\item Isoler la fonction intérieure $u(x)$ et calculer $\displaystyle\lim_{x\to a}u(x)=b$.
		\item Calculer $\displaystyle\lim_{t\to b}g(t)$.
		\item Conclure.
	\end{enumerate}
\end{methode}

\begin{exemple}
	$\displaystyle\lim_{x\to 0^+}\e^{1/x}$ : $u(x)=\dfrac{1}{x}\to+\infty$,
	puis $\e^{t}\to+\infty$. Donc la limite est $+\infty$.
\end{exemple}

\begin{exemple}
	$\displaystyle\lim_{x\to+\infty}\ln\!\left(1+\frac{1}{x}\right)$ :
	$u(x)=1+\dfrac{1}{x}\to 1$, puis $\ln t\to 0$. Donc la limite est $0$.
\end{exemple}

\begin{exemple}
	$\displaystyle\lim_{x\to+\infty}\e^{-x^2}$ : $u(x)=-x^2\to-\infty$,
	puis $\e^{t}\to 0$. Donc la limite est $0$.
\end{exemple}


% ================================================================
\section{Théorèmes de comparaison et des gendarmes}
% ================================================================

\begin{theoreme}[Comparaison]
	Au voisinage de $a$ (fini ou infini), si $f(x)\leq g(x)$ :
	\begin{itemize}[label=\textbullet]
		\item $\displaystyle\lim_{x\to a}f(x)=+\infty
		\;\Longrightarrow\;\lim_{x\to a}g(x)=+\infty$ ;
		\item $\displaystyle\lim_{x\to a}g(x)=-\infty
		\;\Longrightarrow\;\lim_{x\to a}f(x)=-\infty$.
	\end{itemize}
\end{theoreme}

\begin{theoreme}[Théorème des gendarmes]
	Au voisinage de $a$, si $\;g(x)\leq f(x)\leq h(x)\;$ et si
	$\displaystyle\lim_{x\to a}g(x)=\lim_{x\to a}h(x)=\ell$, alors
	$\displaystyle\lim_{x\to a}f(x)=\ell$.
\end{theoreme}

\begin{astucebac}
	Deux inégalités valables sur $\R$, très utiles pour comparer :
	\[
	\forall x\in\R,\quad \e^x\geq x+1
	\qquad\text{et}\qquad
	\forall x>0,\quad \ln x\leq x-1 .
	\]
\end{astucebac}

\begin{exemple}
	Montrer que $\displaystyle\lim_{x\to+\infty}\frac{\ln x}{x}=0$.

	Pour $x>1$, on applique $\ln t\leq t-1$ à $t=\sqrt{x}$ :
	$\ln\sqrt{x}\leq\sqrt{x}-1$, soit $\dfrac{1}{2}\ln x\leq\sqrt{x}$.
	Donc, pour $x>1$,
	\[
	0<\frac{\ln x}{x}\leq\frac{2\sqrt{x}}{x}=\frac{2}{\sqrt{x}}
	\xrightarrow[x\to+\infty]{}0 .
	\]
	Par le théorème des gendarmes, $\displaystyle\lim_{x\to+\infty}\frac{\ln x}{x}=0$.
\end{exemple}

\begin{exemple}
	Montrer que $\displaystyle\lim_{x\to+\infty}\bigl(\e^x-x\bigr)=+\infty$.

	Pour tout $x$, $\e^x\geq x+1$, donc $\e^x-x\geq 1$… ce qui ne suffit pas.
	On applique plutôt $\e^t\geq t+1$ à $t=\dfrac{x}{2}$ :
	$\e^{x/2}\geq\dfrac{x}{2}+1\geq\dfrac{x}{2}$, d'où $\e^{x}\geq\dfrac{x^2}{4}$
	pour $x\geq 0$. Alors
	\[
	\e^x-x\geq\frac{x^2}{4}-x=x\!\left(\frac{x}{4}-1\right)
	\xrightarrow[x\to+\infty]{}+\infty ,
	\]
	et par comparaison $\displaystyle\lim_{x\to+\infty}\bigl(\e^x-x\bigr)=+\infty$.
\end{exemple}


% ================================================================
\section{Continuité}
% ================================================================

\begin{definition}
	Soit $f$ définie sur un intervalle $I$ et $a\in I$.
	$f$ est \textbf{continue en $a$} lorsque
	\[
	\lim_{x\to a}f(x)=f(a)
	\]
	(la limite existe, est finie et vaut $f(a)$).
	$f$ est \textbf{continue sur $I$} si elle est continue en chaque point de $I$.
\end{definition}

\begin{propriete}{Fonctions continues usuelles}{prop_continues_usuelles}
	Sont continues sur leur ensemble de définition :
	\begin{itemize}[label=\textbullet]
		\item les fonctions polynômes et rationnelles ;
		\item $x\mapsto\sqrt{x}$, $x\mapsto|x|$ ;
		\item $x\mapsto\ln x$ sur $]0\,;+\infty[$ et $x\mapsto\e^x$ sur $\R$ ;
		\item toute somme, tout produit, tout quotient (dénominateur non nul)
		et toute \textbf{composée} de fonctions continues.
	\end{itemize}
	En particulier, toute fonction obtenue par opérations et compositions à
	partir de $\ln$ et $\exp$ est continue sur son ensemble de définition.
\end{propriete}

\begin{methode}
	Pour étudier la continuité de $f$ en $a$ (souvent une fonction définie par
	morceaux) :
	\begin{enumerate}
		\item calculer $\displaystyle\lim_{x\to a^-}f(x)$ et
		$\displaystyle\lim_{x\to a^+}f(x)$ ;
		\item les comparer à $f(a)$ ;
		\item conclure : $f$ est continue en $a$ ssi les trois valeurs coïncident.
	\end{enumerate}
\end{methode}

\begin{exemple}
	Soit
	\[
	f(x)=
	\begin{cases}
		\dfrac{\e^x-1}{x} & \text{si } x\neq 0,\\[1.2em]
		1 & \text{si } x=0.
	\end{cases}
	\]
	Sur $\R^*$, $f$ est continue (quotient de fonctions continues).
	En $0$ : $\displaystyle\lim_{x\to 0}\frac{\e^x-1}{x}=1=f(0)$.
	Donc $f$ est continue sur $\R$.
\end{exemple}

\subsection{Prolongement par continuité}

\begin{definition}
	Si $f$ est définie sur $I\setminus\{a\}$ et si
	$\displaystyle\lim_{x\to a}f(x)=\ell\in\R$, on \textbf{prolonge $f$ par
	continuité} en $a$ en posant $f(a)=\ell$. La fonction obtenue est continue
	en $a$.
\end{definition}

\begin{exemple}
	$f(x)=x\ln x$ est définie sur $]0\,;+\infty[$. Comme
	$\displaystyle\lim_{x\to 0^+}x\ln x=0$, on prolonge $f$ par continuité en $0$
	en posant $f(0)=0$.
\end{exemple}

\begin{exemple}
	$g(x)=\dfrac{\ln x}{x-1}$ est définie sur $]0\,;1[\cup\,]1\,;+\infty[$.
	Comme $\displaystyle\lim_{x\to 1}\frac{\ln x}{x-1}=1$, on prolonge $g$ par
	continuité en $1$ en posant $g(1)=1$.
\end{exemple}


% ================================================================
\section{Théorème des valeurs intermédiaires}
% ================================================================

\begin{theoreme}[TVI]
	Soit $f$ \textbf{continue} sur $[a\,;b]$ et $k$ un réel compris entre $f(a)$
	et $f(b)$. Alors il existe \textbf{au moins un} réel $c\in[a\,;b]$ tel que
	$f(c)=k$.
\end{theoreme}

\begin{theoreme}[Corollaire — cas strictement monotone]
	Si de plus $f$ est \textbf{strictement monotone} sur $[a\,;b]$, alors le réel
	$c$ est \textbf{unique}.

	Ce résultat s'étend à un intervalle ouvert ou non borné, en remplaçant
	$f(a)$ et $f(b)$ par les limites de $f$ aux bornes.
\end{theoreme}

\begin{propriete}{Image d'un intervalle}{prop_image_intervalle}
	L'image d'un intervalle par une fonction \textbf{continue} est un intervalle.
	L'image d'un \textbf{segment} $[a\,;b]$ par une fonction continue est un
	segment $[m\,;M]$ ($f$ atteint un minimum et un maximum).
\end{propriete}

\begin{attention}
	Le TVI donne l'\textbf{existence} d'une solution, pas sa valeur. Pour une
	équation $f(x)=0$, il suffit de trouver $a$, $b$ avec
	$f(a)\cdot f(b)<0$ ($f$ continue).
\end{attention}

\begin{methode}
	Pour montrer que $f(x)=k$ a une unique solution sur $I$ :
	\begin{enumerate}
		\item $f$ est continue sur $I$ ;
		\item $f$ est strictement monotone sur $I$ (variations, ou somme/composée
		de fonctions monotones) ;
		\item $k$ est compris entre les valeurs (ou limites) de $f$ aux bornes de $I$ ;
		\item conclure par le corollaire ; encadrer $c$ par balayage si demandé.
	\end{enumerate}
\end{methode}

\begin{center}
\begin{tikzpicture}[>=Stealth, scale=1.05]
	\draw[->] (-0.3,0) -- (4.7,0) node[right,font=\small] {$x$};
	\draw[->] (0,-0.4) -- (0,3.6) node[above,font=\small] {$y$};
	\draw[thick,domain=0.35:4.2,samples=80] plot (\x,{0.28*exp(0.62*\x)});
	\draw[dashed] (0.8,0) node[below,font=\small] {$a$} -- (0.8,{0.28*exp(0.62*0.8)});
	\draw[dashed] (3.8,0) node[below,font=\small] {$b$} -- (3.8,{0.28*exp(0.62*3.8)});
	\draw[dashed,gray] (0,1.7) node[left,font=\small] {$k$} -- (2.53,1.7);
	\draw[dashed,gray] (2.53,0) node[below,font=\small] {$c$} -- (2.53,1.7);
	\filldraw[red] (2.53,1.7) circle (2pt);
	\filldraw (0.8,{0.28*exp(0.62*0.8)}) circle (1.4pt)
	  node[left,font=\tiny] {$f(a)$};
	\filldraw (3.8,{0.28*exp(0.62*3.8)}) circle (1.4pt)
	  node[above left,font=\tiny] {$f(b)$};
	\node[below left,font=\small] at (0,0) {$O$};
\end{tikzpicture}
\end{center}

\begin{exemple}
	Montrer que l'équation $\e^x+x-2=0$ admet une unique solution $\alpha$
	dans $\R$, puis que $0<\alpha<1$.

	$f:x\mapsto\e^x+x-2$ est continue et \textbf{strictement croissante} sur $\R$
	(somme de deux fonctions strictement croissantes).
	$\displaystyle\lim_{x\to-\infty}f(x)=-\infty$ et
	$\displaystyle\lim_{x\to+\infty}f(x)=+\infty$, donc $0$ est atteint une seule
	fois : il existe un unique $\alpha\in\R$ tel que $f(\alpha)=0$.

	$f(0)=1+0-2=-1<0$ et $f(1)=\e-1>0$, donc $\alpha\in\,]0\,;1[$.
\end{exemple}


% ================================================================
\section{Limites et asymptotes}
% ================================================================

\begin{definition}
	Soit $\mathcal C$ la courbe de $f$.
	\begin{itemize}[label=\textbullet]
		\item \textbf{Asymptote horizontale} $y=\ell$ si
		$\displaystyle\lim_{x\to\pm\infty}f(x)=\ell$.
		\item \textbf{Asymptote verticale} $x=a$ si
		$\displaystyle\lim_{x\to a}f(x)=\pm\infty$.
		\item \textbf{Asymptote oblique} $y=mx+p$ si
		$\displaystyle\lim_{x\to\pm\infty}\bigl[f(x)-(mx+p)\bigr]=0$.
	\end{itemize}
\end{definition}

\begin{exemple}
	$f(x)=\dfrac{\e^x}{\e^x+1}$ sur $\R$.
	\begin{itemize}[label=\textbullet]
		\item En $-\infty$ : $\e^x\to 0$, donc $f(x)\to 0$ : asymptote $y=0$.
		\item En $+\infty$ : $f(x)=\dfrac{1}{1+\e^{-x}}\to 1$ : asymptote $y=1$.
	\end{itemize}
\end{exemple}

\begin{exemple}
	$f(x)=x+\e^{-x}$ sur $\R$. En $+\infty$, $f(x)-x=\e^{-x}\to 0$ :
	la droite $y=x$ est asymptote oblique, et $\mathcal C$ est \textbf{au-dessus}
	de cette droite (car $\e^{-x}>0$).
\end{exemple}

\begin{remarque}
	$x\mapsto\ln x$ et $x\mapsto\e^x$ n'ont pas d'asymptote oblique : en
	$+\infty$, $\dfrac{\ln x}{x}\to 0$ (branche parabolique de direction
	$(Ox)$) et $\dfrac{\e^x}{x}\to+\infty$ (branche parabolique de direction
	$(Oy)$).
\end{remarque}


% ================================================================
\section{Méthodes et exercices résolus}
% ================================================================

\begin{methode}
	\textbf{Choisir la bonne technique de levée d'indétermination.}
	\begin{center}
	\renewcommand{\arraystretch}{1.4}
	\begin{tabular}{>{\raggedright\arraybackslash}p{3.7cm}
	                >{\raggedright\arraybackslash}p{4.4cm}}
	\hline
	\textbf{Situation} & \textbf{Réflexe} \\
	\hline
	$\e^x$, $x^n$, $\ln x$ mélangés en $\pm\infty$
	& mettre en facteur le terme dominant \\
	quotient $\dfrac{\infty}{\infty}$
	& diviser haut et bas par le terme dominant \\
	$\dfrac{\e^{\square}-1}{\square}$, $\dfrac{\ln(1+\square)}{\square}$, $\dfrac{\ln x}{x-1}$
	& faire apparaître la limite usuelle ($=1$) \\
	expression en $\dfrac{1}{x}$ ou $\sqrt{x}$
	& changement de variable \\
	$u(x)^{v(x)}$
	& passer à $\e^{v(x)\ln u(x)}$ \\
	terme borné / oscillant, ou inégalité connue
	& comparaison / gendarmes \\
	\hline
	\end{tabular}
	\end{center}
\end{methode}

\exerciceresolu
\textbf{Six limites}

Calculer :
\[
\text{a) }\lim_{x\to+\infty}\frac{\e^x-3x}{\e^x+1}
\qquad
\text{b) }\lim_{x\to+\infty}\frac{x^3-\ln x}{x^2}
\qquad
\text{c) }\lim_{x\to 0}\frac{\ln(1+5x)}{x}
\]
\[
\text{d) }\lim_{x\to-\infty}x\,\e^{x}
\qquad
\text{e) }\lim_{x\to+\infty}\bigl(\ln(x+1)-\ln x\bigr)
\qquad
\text{f) }\lim_{x\to 0^+}x^2\ln x
\]

\solution

\textbf{a)} On divise par $\e^x$ :
$\dfrac{1-3x\e^{-x}}{1+\e^{-x}}\to\dfrac{1-0}{1+0}=1$.

\textbf{b)} $\dfrac{x^3-\ln x}{x^2}=x-\dfrac{\ln x}{x^2}\to+\infty-0=+\infty$.

\textbf{c)} $\dfrac{\ln(1+5x)}{x}=5\cdot\dfrac{\ln(1+5x)}{5x}\to 5\times 1=5$.

\textbf{d)} Croissance comparée : $\displaystyle\lim_{x\to-\infty}x\e^x=0$.

\textbf{e)} $\ln(x+1)-\ln x=\ln\!\dfrac{x+1}{x}=\ln\!\left(1+\dfrac{1}{x}\right)\to\ln 1=0$.

\textbf{f)} Croissance comparée : $\displaystyle\lim_{x\to 0^+}x^2\ln x=0$.

\bigskip
\exerciceresolu
\textbf{Étude des branches infinies}

Soit $f(x)=(x-1)\,\e^{x}$ définie sur $\R$.
\begin{enumerate}
	\item Déterminer $\displaystyle\lim_{x\to-\infty}f(x)$ et
	$\displaystyle\lim_{x\to+\infty}f(x)$.
	\item Préciser les asymptotes éventuelles.
\end{enumerate}

\solution

\textbf{1.} En $-\infty$ : $f(x)=x\e^x-\e^x$. Or $x\e^x\to 0$ et $\e^x\to 0$,
donc $\displaystyle\lim_{x\to-\infty}f(x)=0$.

En $+\infty$ : $x-1\to+\infty$ et $\e^x\to+\infty$, donc
$\displaystyle\lim_{x\to+\infty}f(x)=+\infty$.

\textbf{2.} En $-\infty$, la droite $y=0$ est asymptote horizontale.
De plus $f(x)<0$ pour $x<1$, donc $\mathcal C$ est \emph{sous} l'axe des
abscisses au voisinage de $-\infty$.
En $+\infty$, $\dfrac{f(x)}{x}=\dfrac{(x-1)\e^x}{x}\to+\infty$ :
pas d'asymptote, branche parabolique de direction $(Oy)$.

\bigskip
\exerciceresolu
\textbf{Prolongement et TVI}

Soit $f(x)=\dfrac{\e^x-1}{x}$ pour $x\neq 0$.
\begin{enumerate}
	\item Montrer que $f$ se prolonge par continuité en $0$.
	\item On admet que le prolongement, encore noté $f$, est strictement
	croissant sur $\R$. Montrer que l'équation $f(x)=2$ admet une unique
	solution $\alpha$, et que $1<\alpha<2$.
\end{enumerate}

\solution

\textbf{1.} $\displaystyle\lim_{x\to 0}\frac{\e^x-1}{x}=1$ (limite usuelle),
donc on pose $f(0)=1$ : $f$ est continue sur $\R$.

\textbf{2.} $f$ est continue et strictement croissante sur $\R$, avec
$\displaystyle\lim_{x\to-\infty}f(x)=0$ (car $\e^x\to 0$ et
$\dfrac{-1}{x}\to 0$) et $\displaystyle\lim_{x\to+\infty}f(x)=+\infty$.
Comme $2\in\,]0\,;+\infty[$, le corollaire du TVI donne une unique solution
$\alpha$.

$f(1)=\e-1\approx 1{,}72<2$ et
$f(2)=\dfrac{\e^2-1}{2}\approx 3{,}19>2$, donc $1<\alpha<2$.

\bigskip
\exerciceresolu
\textbf{Gendarmes}

Montrer que $\displaystyle\lim_{x\to+\infty}\frac{\ln x}{x+\e^x}=0$.

\solution

Pour $x>1$, $\ln x>0$ et $x+\e^x>\e^x>0$, donc
\[
0<\frac{\ln x}{x+\e^x}<\frac{\ln x}{\e^x}
=\frac{\ln x}{x}\times\frac{x}{\e^x}.
\]
Or $\dfrac{\ln x}{x}\to 0$ et $\dfrac{x}{\e^x}\to 0$, donc le majorant tend
vers $0$. Par le théorème des gendarmes, la limite cherchée est $0$.


% ================================================================
\section{Exercices d'entraînement}
% ================================================================

% --------------------------------------------------
\subsection{Échauffement : lire un graphique}
% --------------------------------------------------

\exercice{Lecture graphique \hfill \facile}
La courbe $\mathcal C$ ci-dessous représente une fonction $f$ définie sur
$\R\setminus\{1\}$. Les droites $x=1$ et $y=2$ sont asymptotes à $\mathcal C$.

\begin{center}
\begin{tikzpicture}[>=Stealth, scale=0.82]
	\begin{scope}
	\clip (-3.7,-2.8) rectangle (5.2,5.2);
	\draw[thick,domain=-3.4:0.80,samples=140] plot (\x,{2+1/(\x-1)});
	\draw[thick,domain=1.22:4.6,samples=140]  plot (\x,{2+1/(\x-1)});
	\end{scope}
	\draw[->] (-3.6,0) -- (4.9,0) node[right,font=\small] {$x$};
	\draw[->] (0,-2.6) -- (0,4.9) node[above,font=\small] {$y$};
	\draw[dashed,gray] (1,-2.6) -- (1,4.7);
	\node[font=\scriptsize,gray,anchor=west] at (1.15,-2.3) {$x=1$};
	\draw[dashed,gray] (-3.5,2) -- (4.8,2) node[right,font=\scriptsize,gray] {$y=2$};
	\node[font=\small] at (4.1,3.6) {$\mathcal C$};
\end{tikzpicture}
\end{center}

Par lecture graphique :
\begin{enumerate}
	\item Donner $\displaystyle\lim_{x\to-\infty}f(x)$ et
	$\displaystyle\lim_{x\to+\infty}f(x)$.
	\item Donner $\displaystyle\lim_{x\to 1^-}f(x)$ et
	$\displaystyle\lim_{x\to 1^+}f(x)$.
	\item Sur quels intervalles $f$ est-elle continue ?
	\item $f$ est-elle prolongeable par continuité en $1$ ? Justifier.
\end{enumerate}

\exercice{Continue ou non ? \hfill \facile}
Pour chacune des fonctions suivantes, étudier la continuité en $x_0=0$.
\begin{enumerate}
	\item $f(x)=\e^x$ si $x\leq 0$ et $f(x)=x+1$ si $x>0$.
	\item $g(x)=\dfrac{\e^{x}-1}{x}$ si $x\neq 0$ et $g(0)=0$.
	\item $h(x)=x\ln|x|$ si $x\neq 0$ et $h(0)=0$.
\end{enumerate}

\exercice{Vrai ou faux \hfill \facile}
Répondre par \emph{vrai} ou \emph{faux} en justifiant brièvement.
\begin{enumerate}
	\item $\displaystyle\lim_{x\to+\infty}\frac{\e^x}{x^{100}}=+\infty$.
	\item $\displaystyle\lim_{x\to+\infty}\frac{\ln x}{x}=1$.
	\item $\displaystyle\lim_{x\to 0^+}x\ln x=-\infty$.
	\item Une fonction continue sur $[a\,;b]$ y est bornée et atteint ses bornes.
	\item Si $f(a)\cdot f(b)<0$ avec $f$ continue sur $[a\,;b]$, alors $f$
	s'annule au moins une fois sur $]a\,;b[$.
	\item Si $f$ est continue sur $[a\,;b]$ et $f(a)\cdot f(b)>0$, alors $f$ ne
	s'annule pas sur $[a\,;b]$.
	\item $\displaystyle\lim_{x\to 0}\frac{\ln(1+x)}{x}=\lim_{x\to 0}\frac{\e^x-1}{x}$.
	\item Si $\displaystyle\lim_{x\to a}f(x)=f(a)$, alors $f$ est continue en $a$.
\end{enumerate}

% --------------------------------------------------
\subsection{Batteries de calcul de limites}
% --------------------------------------------------

\begin{astucebac}
	Ces listes sont à traiter \textbf{sans rédaction longue} : identifier la
	forme, nommer la technique, appliquer une limite usuelle. Les réponses sont
	regroupées en fin de sous-section.
\end{astucebac}

\exercice{Termes dominants : $\infty-\infty$ et $\frac{\infty}{\infty}$ \hfill \facile}
\begin{multicols}{2}
\begin{enumerate}
	\item $\displaystyle\lim_{x\to+\infty}\frac{\e^x+3}{\e^x-1}$
	\item $\displaystyle\lim_{x\to+\infty}\frac{2\e^x-x}{\e^x+1}$
	\item $\displaystyle\lim_{x\to-\infty}\frac{\e^x+2}{\e^x-3}$
	\item $\displaystyle\lim_{x\to+\infty}\frac{\e^{2x}-\e^x}{\e^{2x}+1}$
	\item $\displaystyle\lim_{x\to+\infty}\frac{\ln x-1}{\ln x+x}$
	\item $\displaystyle\lim_{x\to+\infty}\bigl(\e^x-x^2\bigr)$
	\item $\displaystyle\lim_{x\to+\infty}\bigl(x-\ln x\bigr)$
	\item $\displaystyle\lim_{x\to 0^+}\Bigl(\ln x+\frac{1}{x}\Bigr)$
	\item $\displaystyle\lim_{x\to+\infty}\bigl(\ln x-\ln(x+1)\bigr)$
	\item $\displaystyle\lim_{x\to+\infty}x\bigl(\ln(x+1)-\ln x\bigr)$
\end{enumerate}
\end{multicols}

\exercice{Croissances comparées \hfill \facile}
\begin{multicols}{3}
\begin{enumerate}
	\item $\displaystyle\lim_{x\to+\infty}\frac{\ln x}{x}$
	\item $\displaystyle\lim_{x\to+\infty}\frac{\ln x}{x^2}$
	\item $\displaystyle\lim_{x\to 0^+}x\ln x$
	\item $\displaystyle\lim_{x\to 0^+}x^2\ln x$
	\item $\displaystyle\lim_{x\to 0^+}\sqrt{x}\,\ln x$
	\item $\displaystyle\lim_{x\to+\infty}\frac{\e^x}{x^3}$
	\item $\displaystyle\lim_{x\to-\infty}x\,\e^x$
	\item $\displaystyle\lim_{x\to-\infty}x^2\,\e^x$
	\item $\displaystyle\lim_{x\to+\infty}x^5\,\e^{-x}$
	\item $\displaystyle\lim_{x\to+\infty}\frac{(\ln x)^2}{x}$
	\item $\displaystyle\lim_{x\to+\infty}\frac{x+\ln x}{\e^x}$
	\item $\displaystyle\lim_{x\to+\infty}x\,\e^{-x}$
\end{enumerate}
\end{multicols}

\exercice{Limites usuelles (taux de variation) \hfill \moyen}
\begin{multicols}{2}
\begin{enumerate}
	\item $\displaystyle\lim_{x\to 0}\frac{\e^x-1}{x}$
	\item $\displaystyle\lim_{x\to 0}\frac{\e^{2x}-1}{x}$
	\item $\displaystyle\lim_{x\to 0}\frac{\e^{-x}-1}{x}$
	\item $\displaystyle\lim_{x\to 0}\frac{\e^{x}-1}{3x}$
	\item $\displaystyle\lim_{x\to 0}\frac{\ln(1+3x)}{x}$
	\item $\displaystyle\lim_{x\to 0}\frac{\ln(1-x)}{x}$
	\item $\displaystyle\lim_{x\to 1}\frac{\ln x}{x-1}$
	\item $\displaystyle\lim_{x\to 1}\frac{\ln x}{x^2-1}$
	\item $\displaystyle\lim_{x\to 0}\frac{\e^{3x}-1}{\e^{x}-1}$
	\item $\displaystyle\lim_{x\to 0}\frac{x}{\e^x-1}$
	\item $\displaystyle\lim_{x\to 0}\frac{\e^x-1}{\ln(1+x)}$
	\item $\displaystyle\lim_{x\to 2}\frac{\ln(x-1)}{x-2}$
\end{enumerate}
\end{multicols}

\exercice{Changement de variable \hfill \moyen}
\begin{multicols}{2}
\begin{enumerate}
	\item $\displaystyle\lim_{x\to+\infty}x\bigl(\e^{1/x}-1\bigr)$
	\item $\displaystyle\lim_{x\to+\infty}x\ln\!\Bigl(1+\frac{1}{x}\Bigr)$
	\item $\displaystyle\lim_{x\to+\infty}\frac{\ln x}{\sqrt{x}-1}$
	\item $\displaystyle\lim_{x\to+\infty}\frac{\ln x}{x^{1/3}}$
	\item $\displaystyle\lim_{x\to 0^+}x^{x}$
	\item $\displaystyle\lim_{x\to+\infty}\Bigl(1+\frac{1}{x}\Bigr)^{x}$
	\item $\displaystyle\lim_{x\to+\infty}x^{1/x}$
	\item $\displaystyle\lim_{x\to 0^+}\e^{-1/x}\cdot\frac{1}{x}$
\end{enumerate}
\end{multicols}

\exercice{Composées et gendarmes \hfill \moyen}
\begin{multicols}{2}
\begin{enumerate}
	\item $\displaystyle\lim_{x\to+\infty}\e^{-x^2}$
	\item $\displaystyle\lim_{x\to 0^+}\e^{-1/x}$
	\item $\displaystyle\lim_{x\to 0^+}\e^{1/x}$
	\item $\displaystyle\lim_{x\to+\infty}\ln\bigl(1+\e^{-x}\bigr)$
	\item $\displaystyle\lim_{x\to+\infty}\bigl(\ln(\e^x+1)-x\bigr)$
	\item $\displaystyle\lim_{x\to+\infty}\frac{\ln(\e^x+x)}{x}$
	\item $\displaystyle\lim_{x\to+\infty}\frac{\ln x}{x+\e^x}$
	\item $\displaystyle\lim_{x\to-\infty}\e^{x}\ln(1-x)$
\end{enumerate}
\end{multicols}

\medskip
\begin{tcolorbox}[enhanced,breakable,colback=gray!5,colframe=black!55,
	boxrule=0.4pt,arc=1mm,left=3mm,right=3mm,top=2mm,bottom=2mm,
	fonttitle=\bfseries\sffamily,title={Réponses}]
{\footnotesize
\textbf{Termes dominants.}
1) $1$ \; 2) $2$ \; 3) $-\tfrac{2}{3}$ \; 4) $1$ \; 5) $0$ \; 6) $+\infty$ \;
7) $+\infty$ \; 8) $+\infty$ \; 9) $0$ \; 10) $1$.

\smallskip
\textbf{Croissances comparées.}
1) $0$ \; 2) $0$ \; 3) $0$ \; 4) $0$ \; 5) $0$ \; 6) $+\infty$ \; 7) $0$ \;
8) $0$ \; 9) $0$ \; 10) $0$ \; 11) $0$ \; 12) $0$.

\smallskip
\textbf{Limites usuelles.}
1) $1$ \; 2) $2$ \; 3) $-1$ \; 4) $\tfrac{1}{3}$ \; 5) $3$ \; 6) $-1$ \;
7) $1$ \; 8) $\tfrac{1}{2}$ \; 9) $3$ \; 10) $1$ \; 11) $1$ \; 12) $1$.

\smallskip
\textbf{Changement de variable.}
1) $1$ \; 2) $1$ \; 3) $0$ \; 4) $0$ \; 5) $1$ \; 6) $\e$ \; 7) $1$ \; 8) $0$.

\smallskip
\textbf{Composées et gendarmes.}
1) $0$ \; 2) $0$ \; 3) $+\infty$ \; 4) $0$ \; 5) $0$ \; 6) $1$ \; 7) $0$ \;
8) $0$.
}
\end{tcolorbox}

% --------------------------------------------------
\subsection{Continuité, prolongement et TVI}
% --------------------------------------------------

\exercice{Prolongements \hfill \facile}
Étudier si les fonctions suivantes, définies sur $\R^*$ ou $]0\,;+\infty[$,
se prolongent par continuité en $0$ (resp.\ en $1$).
\begin{multicols}{2}
\begin{enumerate}
	\item $f(x)=\dfrac{\e^{2x}-1}{x}$
	\item $g(x)=\dfrac{\ln(1+x)}{x}$
	\item $h(x)=x\ln x$
	\item $k(x)=\e^{-1/x^2}$ (en $0$)
	\item $m(x)=\dfrac{x-1}{\ln x}$ (en $1$)
	\item $p(x)=\dfrac{\e^x-1}{x^2}$
\end{enumerate}
\end{multicols}

\exercice{Raccord d'une fonction par morceaux \hfill \moyen}
Soit $a$ un réel et
\[
f(x)=
\begin{cases}
	a+\ln(1+x) & \text{si } -1<x\leq 0,\\[0.4em]
	\dfrac{\e^{x}-1}{x} & \text{si } x>0.
\end{cases}
\]
\begin{enumerate}
	\item Calculer $\displaystyle\lim_{x\to 0^+}f(x)$.
	\item Déterminer $a$ pour que $f$ soit continue en $0$.
	\item Pour cette valeur de $a$, $f$ est-elle continue sur $]-1\,;+\infty[$ ?
\end{enumerate}

\exercice{TVI et encadrement \hfill \moyen}
Soit $f(x)=x+\ln x$ sur $]0\,;+\infty[$.
\begin{enumerate}
	\item Déterminer $\displaystyle\lim_{x\to 0^+}f(x)$ et
	$\displaystyle\lim_{x\to+\infty}f(x)$.
	\item Montrer que $f$ réalise une bijection de $]0\,;+\infty[$ sur $\R$.
	\item En déduire que l'équation $f(x)=0$ admet une unique solution $\alpha$.
	\item Montrer que $\tfrac{1}{2}<\alpha<1$, puis que $\alpha=\e^{-\alpha}$.
\end{enumerate}

\exercice{Deux asymptotes horizontales \hfill \moyen}
Soit $f(x)=\dfrac{2\e^x+1}{\e^x+1}$ sur $\R$.
\begin{enumerate}
	\item Déterminer $\displaystyle\lim_{x\to-\infty}f(x)$ et
	$\displaystyle\lim_{x\to+\infty}f(x)$. Interpréter graphiquement.
	\item Montrer que pour tout $x$, $1<f(x)<2$.
	\item Montrer que $f(x)=\dfrac{3}{2}$ admet une unique solution et la
	calculer exactement.
\end{enumerate}

\exercice{Une inégalité pour encadrer \hfill \difficile}
On admet que pour tout $x>0$, $\;\dfrac{x}{x+1}\leq\ln(1+x)\leq x$.
\begin{enumerate}
	\item Retrouver $\displaystyle\lim_{x\to 0^+}\frac{\ln(1+x)}{x}$.
	\item En déduire $\displaystyle\lim_{x\to+\infty}x\ln\!\Bigl(1+\frac{1}{x}\Bigr)$.
	\item Montrer que la suite $u_n=\Bigl(1+\dfrac{1}{n}\Bigr)^{n}$ converge et
	donner sa limite.
\end{enumerate}

% --------------------------------------------------
\subsection{Sujets types Bac}
% --------------------------------------------------

\exercice{Sujet type Bac — fonction par morceaux et branches infinies \hfill \difficile}
On considère la fonction $f$ définie sur $\R$ par
\[
f(x)=
\begin{cases}
	\dfrac{\e^{x}-1}{x} & \text{si } x<0,\\[1em]
	1 & \text{si } x=0,\\[0.4em]
	x\,\e^{-x}+1 & \text{si } x>0,
\end{cases}
\]
et l'on note $\mathcal C$ sa courbe représentative.

\textbf{Partie A — Continuité et limites}
\begin{enumerate}
	\item Montrer que $f$ est continue en $0$.
	\item Montrer que la droite $y=0$ est asymptote à $\mathcal C$ en $-\infty$.
	\item Montrer que la droite $y=1$ est asymptote à $\mathcal C$ en $+\infty$.
	\item Montrer que, pour tout $x>0$, $f(x)>1$ ; préciser la position de
	$\mathcal C$ par rapport à son asymptote en $+\infty$.
\end{enumerate}

\textbf{Partie B — Une équation}
\begin{enumerate}
	\setcounter{enumi}{4}
	\item On admet que $f$ est continue et strictement croissante sur
	$]-\infty\,;0[$. Déterminer $\displaystyle\lim_{x\to-\infty}f(x)$ et
	$\displaystyle\lim_{x\to 0^-}f(x)$.
	\item En déduire que l'équation $f(x)=\dfrac{1}{2}$ admet une unique solution
	$\alpha$ dans $]-\infty\,;0[$.
	\item Montrer que $-2<\alpha<-1$.
	\item Montrer que $\alpha$ vérifie $\e^{\alpha}=1+\dfrac{\alpha}{2}$.
\end{enumerate}

\exercice{Sujet type Bac — la fonction $x\mapsto\dfrac{\ln x}{x}$ \hfill \difficile}
Soit $f$ la fonction définie sur $]0\,;+\infty[$ par $f(x)=\dfrac{\ln x}{x}$,
de courbe $\mathcal C$.

\textbf{Partie A — Limites et variations}
\begin{enumerate}
	\item Déterminer $\displaystyle\lim_{x\to 0^+}f(x)$ ; en déduire une
	asymptote.
	\item Déterminer $\displaystyle\lim_{x\to+\infty}f(x)$ ; en déduire une
	asymptote et la nature de la branche infinie.
	\item On admet que $f$ est strictement croissante sur $]0\,;\e]$ et
	strictement décroissante sur $[\e\,;+\infty[$, avec $f(\e)=\dfrac{1}{\e}$.
	Dresser le tableau de variations.
\end{enumerate}

\textbf{Partie B — Nombre de solutions de $f(x)=m$}
\begin{enumerate}
	\setcounter{enumi}{3}
	\item Déterminer $f\bigl(]0\,;\e]\bigr)$ et $f\bigl([\e\,;+\infty[\bigr)$.
	\item En déduire, suivant les valeurs du réel $m$, le nombre de solutions de
	l'équation $\ln x=mx$.
	\item Montrer que l'équation $f(x)=\dfrac{1}{4}$ admet exactement deux
	solutions $x_1\in\,]1\,;\e[$ et $x_2\in\,]\e\,;+\infty[$, et donner un
	encadrement d'amplitude $1$ de chacune.
\end{enumerate}

\exercice{Sujet type Bac — suite définie implicitement \hfill \difficile}
Pour tout entier $n\geq 1$, on considère l'équation
\[
(E_n):\qquad \e^{x}+x=n .
\]

\textbf{Partie A — Existence de $\alpha_n$}
\begin{enumerate}
	\item Soit $g_n(x)=\e^{x}+x-n$. Montrer que $g_n$ est continue et
	strictement croissante sur $\R$.
	\item Déterminer $\displaystyle\lim_{x\to-\infty}g_n(x)$ et
	$\displaystyle\lim_{x\to+\infty}g_n(x)$.
	\item En déduire que $(E_n)$ admet une unique solution $\alpha_n$ dans $\R$.
	\item Montrer que $0<\alpha_n<\ln n$ pour $n\geq 2$.
\end{enumerate}

\textbf{Partie B — Étude de la suite $(\alpha_n)$}
\begin{enumerate}
	\setcounter{enumi}{4}
	\item Montrer que $g_{n+1}(\alpha_n)<0$. En déduire que $(\alpha_n)$ est
	strictement croissante.
	\item Montrer que $\displaystyle\lim_{n\to+\infty}\alpha_n=+\infty$.
	\item À partir de $\e^{\alpha_n}=n-\alpha_n$, montrer que
	$\alpha_n=\ln n+\ln\!\Bigl(1-\dfrac{\alpha_n}{n}\Bigr)$, puis que
	$\displaystyle\lim_{n\to+\infty}\bigl(\alpha_n-\ln n\bigr)=0$.
\end{enumerate}

\exercice{Sujet type Bac — fonction auxiliaire et TVI \hfill \difficile}
Soit $f$ la fonction définie sur $\R$ par $f(x)=x\,\e^{x}-1$, de courbe
$\mathcal C$.

\textbf{Partie A — Limites et signe}
\begin{enumerate}
	\item Déterminer $\displaystyle\lim_{x\to-\infty}f(x)$ ; en déduire une
	asymptote.
	\item Déterminer $\displaystyle\lim_{x\to+\infty}f(x)$.
	\item On admet que $f$ est continue et strictement croissante sur
	$[0\,;+\infty[$. Montrer que l'équation $f(x)=0$ admet une unique solution
	$\alpha$ dans $[0\,;+\infty[$, et qu'elle n'en admet aucune sur
	$]-\infty\,;0]$.
	\item Montrer que $0{,}5<\alpha<0{,}6$.
	\item Montrer que $\alpha=\e^{-\alpha}$ et que $\alpha=-\ln\alpha$.
\end{enumerate}

\textbf{Partie B — Une suite d'approximations}
\begin{enumerate}
	\setcounter{enumi}{5}
	\item On pose $u_0=\dfrac{1}{2}$ et, pour tout $n$, $u_{n+1}=\e^{-u_n}$.
	Calculer $u_1$ et $u_2$ (valeurs approchées à $10^{-3}$).
	\item En admettant que $(u_n)$ converge, déterminer sa limite.
\end{enumerate}

\exercice{Sujet type Bac — symétrie et asymptotes \hfill \difficile}
Soit $f$ la fonction définie sur $\R\setminus\{0\}$ par
$f(x)=\dfrac{\e^{x}}{\e^{x}-1}$, de courbe $\mathcal C$.

\textbf{Partie A — Étude aux bornes}
\begin{enumerate}
	\item Déterminer $\displaystyle\lim_{x\to-\infty}f(x)$ et
	$\displaystyle\lim_{x\to+\infty}f(x)$. En déduire deux asymptotes.
	\item Déterminer $\displaystyle\lim_{x\to 0^-}f(x)$ et
	$\displaystyle\lim_{x\to 0^+}f(x)$. En déduire une asymptote et justifier
	que $f$ \emph{n'est pas} prolongeable par continuité en $0$.
\end{enumerate}

\textbf{Partie B — Propriétés}
\begin{enumerate}
	\setcounter{enumi}{2}
	\item Montrer que $f(x)-1=\dfrac{1}{\e^{x}-1}$ et étudier la position de
	$\mathcal C$ par rapport à la droite $y=1$.
	\item Montrer que, pour tout $x\neq 0$, $f(x)+f(-x)=1$. Interpréter
	géométriquement.
	\item Montrer que l'équation $f(x)=-1$ admet une unique solution que l'on
	déterminera exactement.
\end{enumerate}
